\subsection{Cyclic four-object support at eighteen arrows}

\begin{lemma}[Proved: a one-row source component reduces to its minimal group]
\label{CP2:A-source-minimal-group-collapse}
Let $S$ be a strong source component on at least two objects in
a finite category $C$. If each local minimal row count of $S$
is one, there is an actual homotopy equivalence from a category
with strictly fewer arrows to $C$, replacing $S$ by its finite
minimal group. There is a one-column sink dual.
\end{lemma}
\begin{proof}
Choose a minimal idempotent $e\in S(a,a)$ and adjoin its retract
object to obtain the full one-object extension $C^e$ in the
idempotent completion. Its group corner is $G=eS(a,a)e$.
The added object is a retract of $a$, so
$BC\to BC^e$ is a homotopy equivalence. Let $D$ be the full
complement of $S$ and let $C'$ be the full category on $D$ and
the new object $e$.

For each retained object the comma category of the inclusion
$C'\to C^e$ has its usual terminal identity. For $x\in S$,
there are no arrows from $D$ into $x$, because $S$ is a source
component. The comma category therefore consists of
$C^e(e,x)=S(a,x)e$, with its right $G$ action. The categorical
minimal-ideal lemma makes this action free, and row count one
makes it transitive. There is exactly one morphism between
each ordered pair of its objects, so the comma category is
contractible. Quillen's Theorem~A proves $BC'\simeq BC$.

For $d\in D$, the new outgoing stalk $C^e(e,d)=C(a,d)e$ is a
subset of the old stalk $C(a,d)$. No other outside hom-set grows,
and there are no incoming stalks from $D$. The internal part
$S$ has at least $|G|\,|\operatorname{Ob}S|^2>|G|$ arrows,
whereas its replacement has exactly $|G|$. The total arrow
count strictly decreases. Taking opposites proves the dual.
\end{proof}

\begin{lemma}[Proved: two supported augmentation modules at a small strong pair]
\label{CP2:A-small-pair-augmentation-Tor}
Let $S$ be a strongly connected two-object category with at most
seven arrows. Suppose its categorical minimal ideal $K$ has
exactly one arrow in each ordered hom-set and contains neither
object identity. Adjoin a trivial initial corner $u$ and a trivial
final corner $v$, with every entrance and exit stalk nonempty
and any nonempty direct hom-set $u\to v$. If each of the entrance
and exit augmentation modules is supported on at most one object
of $S$, bounded rational homology of the resulting category
implies that its fourth rational homology is zero.
\end{lemma}
\begin{proof}
Write $R=\mathbb QS$ and let $V,W$ be the left entrance and
right exit augmentation modules. If either is zero, the
middle-bar formula of Lemma~\ref{CP2:A-high-linear-Tor} already
gives the conclusion. Otherwise, at either object the local
minimal idempotent factors through the other object, using
the two minimal cross arrows. It therefore acts by zero on
a module supported at that object. The other minimal arrows
also act by zero, since one of their module spaces is zero.
Thus $K$ annihilates both $V$ and $W$.

Put $A=R/\mathbb QK$. The bimodule description
$\mathbb QK\cong Re\otimes_{\mathbb Q}eR$ in
Lemma~\ref{CP2:A-high-Rees} makes the ideal projective on both
sides. Since $We=0$, the resolution
$0\to\mathbb QK\to R\to A\to0$ and change of rings give
\[
 \operatorname{Tor}_j^R(W,V)\cong
 \operatorname{Tor}_j^A(W,V)\quad(j\ge0).
\]
The two object identities give distinct nonzero orthogonal
idempotents in $A$, so $A$ is nonlocal; its dimension is at
most $7-4=3$. By Lemma~\ref{CP2:A-high-small-Tor}, a nonzero
second Tor group would force nonzero Tor in arbitrarily high
degrees. The exact middle-bar shift identifies these with
homology in degrees two higher. Boundedness therefore forces
$H_4=\operatorname{Tor}_2^A(W,V)=0$.
\end{proof}

\begin{theorem}[Proved: cyclic four-object support through eighteen]
\label{CP2:A-strong-four-eighteen}
A finite category on four objects with cyclic support and at
most eighteen arrows cannot have classifying space of type $\CP^2$.
\end{theorem}
\begin{proof}
Assume a connected counterexample. It has bounded rational
homology and a one-dimensional second cohomology group with
nonzero square. Every actual arrow-reducing homotopy equivalence
would contradict the already proved global lower bound
Theorem~\ref{CP2:A-global-lower-eighteen}; in particular no
object is a retract of another.

A strong component on $t$ objects has its identities outside
the categorical minimal ideal and costs at least $t^2+t$ arrows.
A wholly strong four-object category therefore costs at least
twenty. A proper three-object component costs at least twelve
and its interface at least three, leaving a corner of order
at most three. The strong component has finite-group type:
otherwise its total row and column counts are at least four
and its internal cost is already nineteen. Both blocks have
zero second rational cohomology, and the two-block lemma applies.
If there are two strong pairs, the interface costs at least
four. Each pair then has at most eight internal arrows and
finite-group type, giving the same obstruction.

It remains to consider one strong pair $S$ and two singleton
components. Put $q=|S|\ge6$.

\emph{An extreme pair.} By opposition suppose $S$ is a source.
Let $D$ be its two-object complement and let $h$ count all
interface arrows. Connectedness gives $q\le13$ and $h\ge2$.
If both blocks have zero second rational cohomology, the
two-block lemma finishes.

First suppose $H^2(BD;\mathbb Q)\ne0$. Then $|D|\ge6$:
for connected $D$ this is the second-homology cost bound,
and for disconnected $D$ one corner has order at least five.
Consequently $q\le10$. The internal pair has finite-group
type, since total rows and columns both at least three
would cost at least eleven. If its total row count is two,
Lemma~\ref{CP2:A-source-minimal-group-collapse} is an actual
arrow-reducing equivalence, excluding the case even when its
minimal group is nontrivial.

Otherwise its total column count is two, its minimal group
is trivial, and its total row count is three or four. In
particular $q\ge8$. At $q=8$ it consists exactly of its
six-element minimal band and its two identities. The object
with one local row has all incoming internal arrows fixed by
its minimal idempotent and no units other than its identity.
There are no external entrances, so the one-sided deletion
of Lemma~\ref{CP2:A-high-one-sided} applies. If $h=2$ at
any of the remaining orders, every external stalk is singleton
or empty. The fixed-exit column version of the same lemma
collapses the whole pair and strictly decreases arrows.

The only remaining budget is $q=9$, $|D|=6$, $h=3$.
The complement is connected: two separate components would
require at least four interface arrows. The pair meets only
one object $u$ of $D$, and its other object $v$ points to $u$;
the opposite edge direction would force another interface.
Nonzero second cohomology on a connected six-arrow two-object
category forces corner orders $3,1$ and stalk size two.
Indeed a singleton stalk is a join; if both corners have
orders at most two the semisimple height-one bar has zero
second homology; and a corner of order at least four leaves
only a singleton stalk in this budget.

If the corner at $u$ has order three, all stalks into $u$
from the other three objects have size one or two: the two
from $S$ total three and the one from $v$ has size two.
Lemma~\ref{CP2:A-high-small-binary} applies with the arbitrary
source block $S\sqcup\{v\}$ and kills the square. If the
corner at $u$ is trivial, the full category on $S,u$ is an
ordinary act cone with one or two states at each object,
so Lemma~\ref{CP2:A-high-affine} kills its products. The
full category $D$ also has zero degree-two products by the
two-block lemma, since its corners have orders three and one.
Their nerves cover the whole category and meet at the point
$u$. Fourth-cohomology restriction is injective, so the
global products vanish too.

Now suppose $H^2(BS;\mathbb Q)\ne0$. Then $11\le q\le13$,
the minimal group is trivial, and both total row and column
counts are three. Thus each local row count is at most two,
and Corollary~\ref{CP2:A-baseline-two-row-instances} applies
to ordinary acts over $S$.
If $D$ is disconnected, connectedness requires four interface
stalks, giving $|D|\le3$. Both outside corners have order at
most two. Their exact orbit functors replace them by trivial
leaves. Merging those two vertices is an isomorphism on
cohomology in degrees at least two, respecting products: it
adds one circle to the connected realization. The result is
an ordinary act cone over $S$, whose degree-two products are
zero by the binary projective-cover lemma. Pullback excludes
the original category.

If $D$ is connected and $S$ meets both objects, the budget
forces $q=11$, $|D|=3$, $h=4$. All cross arrows between
$S$ and $D$ are singleton, so this is a categorical join
and has zero positive products. Otherwise $S$ meets only
the final object $u$ of the outside edge $v\to u$. Let $E$
be the full category on $S,u$. The corner at $u$ has order
at most three, since
\[
 |\operatorname{End}(u)|\le18-q-h-2\le3.
\]
The nerves of $E$ and $D$ cover the whole nerve and meet
at this rationally acyclic corner. Fourth restriction is
injective. Also $|D|\le5$, so $H^2(BD;\mathbb Q)=0$.
If the corner at $u$ has order at most two, exact orbits
and the binary projective-cover lemma kill the products
of $BE$. If its order is three, equality in the budget
forces $q=11,h=2,|D|=5$, and $E$ is a singleton-stalk
join with zero products. Restriction therefore kills the
global square in every case. The two extreme branches
exhaust the possibilities, since nonzero second cohomology
on both blocks would require at least $11+6+2=19$ arrows.

\emph{A nonextreme pair.} The condensation is the full chain
$u<S<v$. Let the extreme corner orders be $b,c$, put
$z=|C(u,v)|$, and write the two entrance and exit sizes
as $\alpha_i,\beta_i$ for $i=0,1$. All are positive. Set
\[
 P=\sum_i(\alpha_i-1)+\sum_i(\beta_i-1),\qquad
 \delta=(b-1)+(c-1)+(z-1).
\]
The complete arrow count is
\[
 q+7+P+\delta\le18,\qquad q\le11,
 \qquad P+\delta\le11-q.
\]

For $q\ge8$ and $z\le2$, both endpoint corners have order
at most four. Failure of the binary-stalk test at each end
would require one entrance of size at least three and one
exit of size at least three, costing $P\ge4$. This is
impossible because $P\le3$. Thus one endpoint satisfies
Lemma~\ref{CP2:A-high-small-binary}. For $q\ge8$ and
$z\ge3$, the budget gives $P\le1$ and both endpoint corners
have order at most two. Apply their actual exact-orbit
functors to make them trivial. At least one whole entrance
or exit side originally consisted of singletons, and still
does. Its augmentation module is zero. The middle-bar
formula gives $H_4=0$, again a contradiction.

It remains that $q=6$ or $7$. The minimal group is trivial,
there is one minimal row and column at each object, and
the local minimal idempotents $e_i$ are corner zeros.
Consider the four conditions that $e_i$ fix all incoming
cross arrows, or all outgoing cross arrows. Any one of
these conditions permits the one-sided object deletion;
the zero guarantees its unit condition.

At $q=6$, all internal cross arrows are fixed. If no
condition holds, all four external stalks must be
nonsingleton, costing $P\ge4$. Then $P+\delta\le5$ implies
both endpoint corners and the direct stalk have size at
most two. An entrance and an exit both of size at least
three, in addition to the other two nonsingleton stalks,
would cost $P\ge6$. Hence at least one endpoint has all
three incident stalks binary, and the small-binary lemma
excludes its square.

At $q=7$, the single extra internal arrow can violate at
most two of the four fixed conditions, and a nonsingleton
external stalk can violate only its one incident condition.
Thus $P\le1$ gives a one-sided deletion. Assume $P\ge2$.
If $z\le2$, both endpoint corners have order at most four.
Unless a binary endpoint already applies, one entrance
and one exit have size at least three, so the budget
forces $P=4$, $\delta=0$. Thus $b=c=z=1$, and both the
entrance pair and the exit pair have sizes $(3,1)$ in
some order. Their augmentation modules are each supported
on one object of $S$. Lemma~\ref{CP2:A-small-pair-augmentation-Tor}
gives $H_4=0$ under boundedness.

Finally $q=7$, $z\ge3$, $P\ge2$ forces
$P=2$, $b=c=1$, $z=3$. If one whole side is singleton,
its augmentation module is zero and the middle-bar
formula finishes. Otherwise each side has excess one,
so each augmentation module is one-dimensional and
supported at one object. The same small-pair Tor lemma
applies. This exhausts every nonextreme budget, and
therefore every cyclic four-object category in the
eighteen-arrow range.
\end{proof}
